%! Principal Component Analysis — Unit 7, Lesson 7.3 — Slides (Beamer)
\documentclass[aspectratio=169, 11pt]{beamer}
\usepackage{linalg-beamer}   % provides \burgundyheader{} and \sectionlabel[color]{}
\usetikzlibrary{arrows.meta}
\newcommand{\uu}{\mathbf{u}}
\newcommand{\vv}{\mathbf{v}}
\newcommand{\zero}{\mathbf{0}}
\newcommand{\T}{^{\top}}

\begin{document}

% TITLE SLIDE (CourseName is not defined in beamer, so the name is literal)
{
\setbeamercolor{background canvas}{bg=burgundy}
\begin{frame}[plain]
  \vspace{2.8cm}\hspace{0.55cm}%
  \begin{minipage}{0.9\textwidth}
    {\color{goldacc}\bfseries\small LESSON 7.3}\\[0.5cm]
    {\color{white}\bfseries\LARGE Principal Component Analysis}\\[1.0cm]
    {\color{blushmid}\small Linear Algebra~$\cdot$~Unit 7: The Singular Value Decomposition}
  \end{minipage}
\end{frame}
}

% HOOK
\begin{frame}[t, plain]
  \burgundyheader{One number instead of two?}
  \sectionlabel{Hook}
  \begin{columns}[T]
    \begin{column}{0.52\textwidth}
      Four students, two quiz scores each: $(5,6),(6,5),(3,2),(2,3)$. The scatter leans along a
      diagonal --- do well on one, do well on the other.\\[6pt]
      Lesson~7.2 compressed an image by keeping the biggest singular values. \textbf{PCA does the same to
      data.}
    \end{column}
    \begin{column}{0.44\textwidth}
      \begin{block}{The question}
        Describe each student with a \textbf{single} number instead of two --- along what direction, and how
        much do you lose?
      \end{block}
    \end{column}
  \end{columns}
  \vspace{2pt}
  \begin{block}{Today}
    Center the data, form $A\T A$, and read its top eigenvector: the \textbf{first principal component}.
  \end{block}
\end{frame}

% CENTER
\begin{frame}[t, plain]
  \burgundyheader{Step 0: center the data}
  \sectionlabel[goldacc]{Idea}
  \begin{itemize}
    \setlength{\itemsep}{6pt}
    \item Write the data as a matrix $A$: one \emph{row} per point, one \emph{column} per feature.
    \item \textbf{Center} it --- subtract each column's mean --- so the cloud sits at the origin:
          \[ (5,6),(6,5),(3,2),(2,3)\ \xrightarrow{\ -\,(4,4)\ }\ (1,2),(2,1),(-1,-2),(-2,-1). \]
    \item Now a direction describes how the data \emph{spreads}, not where it sits.
  \end{itemize}
\end{frame}

% PRINCIPAL COMPONENTS
\begin{frame}[t, plain]
  \burgundyheader{The principal components are eigenvectors of $A\T A$}
  \sectionlabel{Skill}
  \begin{columns}[T]
    \begin{column}{0.52\textwidth}
      Same recipe as Lesson~7.1: form the symmetric $A\T A$, take its eigenvectors.
      \[ A\T A=\begin{bsmallmatrix} 10 & 8\\ 8 & 10 \end{bsmallmatrix},\quad \lambda=18,2. \]
      Symmetric $\Rightarrow$ the components are \textbf{perpendicular}:
      \[ \vv_1=\tfrac1{\sqrt2}\begin{bsmallmatrix} 1\\1 \end{bsmallmatrix}\ (\text{PC1}),\quad
         \vv_2=\tfrac1{\sqrt2}\begin{bsmallmatrix} 1\\-1 \end{bsmallmatrix}\ (\text{PC2}). \]
    \end{column}
    \begin{column}{0.46\textwidth}
      \centering
      \begin{tikzpicture}[scale=0.72, font=\scriptsize]
        \draw[linegray, -{Latex[length=1.4mm]}] (-2.3,0)--(2.3,0);
        \draw[linegray, -{Latex[length=1.4mm]}] (0,-2.3)--(0,2.3);
        \foreach \x/\y in {0.8/1.6, 1.6/0.8, -0.8/-1.6, -1.6/-0.8}
          \fill[charcoal] (\x,\y) circle (2pt);
        \draw[burgundy, -{Latex[length=2.2mm]}, very thick] (0,0)--(1.8,1.8);
        \node[burgundy] at (1.35,2.05) {\textbf{PC1}};
        \draw[royalblue, -{Latex[length=2.2mm]}, very thick] (0,0)--(0.8,-0.8);
        \node[royalblue] at (1.45,-0.95) {\textbf{PC2}};
      \end{tikzpicture}
    \end{column}
  \end{columns}
\end{frame}

% READ + INTERPRET
\begin{frame}[t, plain]
  \burgundyheader{How much spread --- and what it means}
  \sectionlabel{Skill}
  \begin{itemize}
    \setlength{\itemsep}{6pt}
    \item The eigenvalue \emph{is} the spread along a direction: $\sigma_1^2=18$, $\sigma_2^2=2$.
    \item \textbf{PC1's share:} $\dfrac{\sigma_1^2}{\sigma_1^2+\sigma_2^2}=\dfrac{18}{20}=\mathbf{90\%}$.
    \item \textbf{Meaning:} PC1 $=\tfrac1{\sqrt2}(1,1)$ is the ``overall score'' axis; PC2 $=\tfrac1{\sqrt2}(1,-1)$
          is ``which quiz was better.'' Overall score is what really separates these students.
  \end{itemize}
  \vspace{2pt}
  \begin{block}{Dimension reduction}
    Summarize each student by their PC1 coordinate alone --- one number instead of two --- and keep $90\%$ of
    the information.
  \end{block}
\end{frame}

% BIG PICTURE + ROAD AHEAD
\begin{frame}[t, plain]
  \burgundyheader{Same SVD, now on data --- and the road ahead}
  \sectionlabel[goldacc]{Payoff}
  \begin{itemize}
    \setlength{\itemsep}{6pt}
    \item Image compression kept the biggest-$\sigma$ \emph{layers}; PCA keeps the biggest-$\sigma$
          \emph{directions}. One rule: \textbf{keep the biggest $\sigma$}.
    \item With $100$ features, the top $2$--$3$ components often carry almost all the spread --- replace $100$
          numbers per point with a handful.
  \end{itemize}
  \vspace{2pt}
  \begin{block}{The road ahead}
    \textbf{7.4} the victory of orthogonality --- why perpendicular directions make the SVD, compression, and
    PCA all work.
  \end{block}
\end{frame}

\end{document}
